% Switch to arabic numerals for main content, starting at page 1 with the Introduction chapter
\pagenumbering{arabic}
\setcounter{page}{1}

% ============================================================
%  CHAPTER 1 — Introduction
%  Typical length: 10–15 pages
%  Cover: background, motivation, problem statement, objectives,
%         and organisation of the thesis.
% ============================================================

% Don't change the section names in this chapter. You can add subsections if needed, but keep the main sections as they are.

\chapter{Introduction}
\label{chap:introduction}

% TODO: Introduce the broad research area in 2–3 paragraphs.
% What is the technology / method / system you are studying?
% Why is it important? What is the current state of practice?

% Example of how to cite a reference inline:
% \citet{key} → Author (Year)   e.g. "Nap and Tuinstra (2023) describe..."
% \citep{key} → (Author, Year)  e.g. "...continuous strain sensing \citep{Nap2023}."

Distributed acoustic sensing (DAS) converts a standard optical fibre into a continuous
array of seismic receivers, providing dense spatial coverage with no downhole electronics
\citep{Willis2022}. [Continue your introductory paragraph here — describe the technology,
its importance, and the current state of practice.]

% Example of how to include a figure:
\begin{figure}[htbp]
  \centering
  \includegraphics[width=\textwidth]{fig1_das_schematic}
  \caption{Operating principle of a DAS system. A laser pulse propagates through the
  fibre and generates Rayleigh backscatter; the phase difference between two points
  separated by the gauge length $GL$ recovers the local axial strain continuously
  along the fibre. Adapted from \citet{Nap2023}.}
  \label{fig:das_schematic}
\end{figure}

% Cross-reference the figure in your text like this:
% "...as illustrated in Figure~\ref{fig:das_schematic}."

[Don't change the section names/headings in this chapter.]

% ---- 1.1 Motivation ----------------------------------------
\section{Motivation}
\label{sec:motivation}

% TODO: Explain why this specific problem matters.
% What gap exists in current practice or understanding?
% What will be lost if this problem is not addressed?

[Write the motivation for your thesis here.]

% ---- 1.2 Problem Statement ---------------------------------
\section{Problem Statement}
\label{sec:problem_statement}

% TODO: Narrow down to the specific research gap you address.
% Be precise about what is missing in the literature.
% An itemised list of gaps works well here.

[State the specific research problem clearly here.]

% ---- 1.3 Objectives ----------------------------------------
\section{Objectives}
\label{sec:objectives}

The objectives of this thesis are as follows.

\begin{enumerate}
  \item Objective 1 — state clearly and specifically.
  \item Objective 2.
  \item Objective 3.
  % Add or remove items as needed — most M.Tech theses have 3–5 objectives.
\end{enumerate}

% ---- 1.4 Organisation of the Thesis ------------------------
\section{Organisation of the Thesis}
\label{sec:organisation}

% TODO: One short paragraph per chapter describing its content.

\textbf{Chapter 2} provides the background study and a literature review.

\textbf{Chapter 3} describes the methodology.

\textbf{Chapter 4} presents the results and discussion.

\textbf{Chapter 5} draws the conclusions and outlines limitations and future work.

\textbf{Appendices} provide supporting material for example - derivations, etc. (if required).


% Uncomment the below lines (if you have used AI for the contents of this chapter) to include an AI use disclosure at the end of the chapter. You can modify the wording as needed, but be transparent about how you used AI tools in writing this chapter.

% \vspace{1.0cm}
% \noindent\rule{\textwidth}{0.4pt}
% \vspace{0.2cm}

% \noindent\textit{Note on AI assistance:} An AI language model was used in writing this chapter as a collaborative tool for [describe how — e.g., literature search, brainstorming, drafting]. All decisions, interpretations, and conclusions are the author's own.
